% Master previous-work review; independent from the existing decks. Style copied from 2026-11-04-Helix.tex.
% Evidence: GreenCache arXiv:2505.23970v3, Secs. 5--6; PRISM
% arXiv:2609.35569v1, Props. 1/3, Sec. 5.3, Appendix A.6, Table 12.
% Conventional GPU serving optimization; no HBM/HBF hardware co-design claims.
\documentclass[aspectratio=169,10pt]{beamer}
\usepackage[english]{babel}
\usepackage[utf8]{inputenc}
\usepackage[T1]{fontenc}
\usepackage{lmodern}
\usepackage{amsmath,amssymb}
\usepackage{booktabs,tabularx,array}
\usepackage{graphicx}
\graphicspath{{./}{figures/}}
\usepackage{xcolor}
\usepackage{tikz}
\usetikzlibrary{arrows.meta,positioning,calc,shapes.geometric}
\usepackage{hyperref}

\usetheme{Madrid}
\usecolortheme{default}
\definecolor{Ink}{HTML}{243748}
\definecolor{Muted}{HTML}{667B8B}
\definecolor{Accent}{HTML}{347E74}
\definecolor{AccentBlue}{HTML}{507DAD}
\definecolor{SoftMint}{HTML}{E6F4EE}
\definecolor{SoftBlue}{HTML}{E9F2F9}
\definecolor{SoftLilac}{HTML}{F1EEF8}
\definecolor{Paper}{HTML}{FBFCFE}
\definecolor{Divider}{HTML}{D8E4EB}

\setbeamercolor{background canvas}{bg=Paper}
\setbeamercolor{normal text}{fg=Ink}
\setbeamercolor{structure}{fg=Accent}
\setbeamercolor{frametitle}{fg=Ink,bg=SoftBlue}
\setbeamercolor{title}{fg=Ink}
\setbeamercolor{subtitle}{fg=Muted}
\setbeamercolor{block title}{fg=Ink,bg=SoftBlue}
\setbeamercolor{block body}{fg=Ink,bg=white}
\setbeamercolor{item}{fg=Accent}
\setbeamercolor{caption name}{fg=Muted}
\setbeamercolor{footline}{fg=Muted,bg=Paper}
\setbeamerfont{frametitle}{size=\large,series=\bfseries}
\setbeamerfont{footline}{size=\tiny}
\setbeamerfont{block title}{series=\bfseries}
\setbeamertemplate{headline}{}
\setbeamertemplate{navigation symbols}{}
\setbeamertemplate{blocks}[rounded][shadow=false]
\setbeamertemplate{itemize items}[circle]
\setbeamertemplate{enumerate items}[default]
\setbeamertemplate{itemize subitem}[circle]
\setbeamertemplate{footline}{}
\setbeamertemplate{frametitle}{%
  \nointerlineskip
  \begin{beamercolorbox}[wd=\paperwidth,ht=1.05cm,dp=.16cm,leftskip=.55cm,rightskip=.40cm]{frametitle}%
    \parbox[b]{\dimexpr\paperwidth-.95cm\relax}{\raggedright\usebeamerfont{frametitle}\fontsize{11}{12}\selectfont\insertframetitle}
  \end{beamercolorbox}}


\setbeamersize{text margin left=.45cm,text margin right=.45cm}
% Two balanced sections share the same geometry on both summary slides.
\newenvironment{summarysection}[1]{%
  \par\noindent{\large\bfseries\color{Accent}#1}\par\vspace{.10cm}
  \begin{itemize}\setlength{\itemsep}{.12cm}\setlength{\parskip}{0pt}
}{\end{itemize}}
\newcommand{\key}[1]{\textbf{#1}}

\newcommand{\categorycover}[3]{%
\begin{frame}[c]{Research Area}
\relax
{\LARGE\bfseries #1\par}\bigskip
{\large\color{Accent}#2\par}\bigskip
{\color{Muted}#3\par}
\end{frame}}
\newcommand{\reviewpoint}[2]{\par{\bfseries\color{Accent}#1}\par\smallskip #2\par\medskip}
\hypersetup{pdftitle={Previous Work Review and Research Directions},
 pdfsubject={Energy-Efficient LLM Serving and HBM/HBF Co-Design}}
\begin{document}
\begin{frame}[c]{Previous Work Review \& Research Directions}
\relax
{\LARGE\bfseries Previous Work Review\par\medskip \& Research Directions\par}
\bigskip
{\large\color{Accent}Energy-Efficient LLM Serving and HBM/HBF Co-Design}\par
\bigskip
\textbf{Research meeting: 4 November 2026}\\
Yu-Jui Liang\quad Advisor: Hsiang-Yun Cheng\\
National Taiwan University / Academia Sinica
\end{frame}
\begin{frame}[c]{Contents}
\renewcommand{\arraystretch}{1.6}
\begin{tabularx}{\textwidth}{@{}Xr@{}}
\textbf{\color{Accent}Research category} & \textbf{\color{Accent}Papers} \\
\midrule
Power-, Energy-, and Carbon-Aware GPU Serving Optimization & 2 \\
HBF System Characterization and KV-Centric Serving & 3 \\
HBM/HBF Architecture and Data-Path Co-Design & 4 \\
HBF Simulation and Power/Thermal Modeling & 3 \\
HBF-Accelerated Retrieval and RAG Systems & 1 \\
\midrule
\textbf{Total} & \textbf{13} \\
\end{tabularx}
\end{frame}
% Category 1; sources are intentionally not shown on slides.
\categorycover{Power-, Energy-, and Carbon-Aware\\GPU Serving Optimization}{How can conventional GPU serving systems optimize energy, carbon, and resource allocation?}{GreenCache; PRISM / Beyond Energy}
% Source: https://arxiv.org/pdf/2505.23970v3
\begin{frame}[t]{Cache Your Prompt When It's Green ---\\Carbon-Aware Caching for Large Language Model Serving}
\vspace{.02cm}
\begin{minipage}[t][3.25cm][t]{\linewidth}
\begin{summarysection}{Motivation}
\item \key{Caching trade-off:} Reusing SSD-based prefix KV caches avoids prefill computation and operational carbon from electricity, but adds embodied carbon from SSD manufacturing.
\item \key{Variable benefits:} Carbon-efficient cache capacity changes with request load and grid carbon intensity. Larger caches are not always greener.
\item \key{Outcome:} Llama-3 70B in France: 15.1\% average carbon reduction vs.\ Full Cache, peaking at 25.3\% in multi-turn serving. Reported latency SLO attainment: $>90\%$.
\end{summarysection}
\end{minipage}
\vspace{.18cm}
\begin{summarysection}{Method}
\item \key{Profile:} Measure latency and power across cache capacities and request loads. Estimate operational and allocated embodied carbon, capturing KV-reuse benefits.
\item \key{Allocate:} Predict load and grid carbon intensity. An ILP selects provisioned SSD cache capacity to minimize carbon subject to TTFT/TPOT SLO-attainment constraints.
\item \key{Replace:} Least Carbon Savings (LCS) scores KV entries using reused tokens, hit frequency, size and age, evicting entries with the lowest carbon-benefit scores.
\end{summarysection}
\end{frame}

% Source: https://arxiv.org/abs/2609.35569
\begin{frame}[t]{Beyond Energy: When Sustainability Dimensions\\Reshape LLM Serving Decisions}
\vspace{.02cm}
\begin{minipage}[t][3.25cm][t]{\linewidth}
\begin{summarysection}{Motivation}
\item \key{Sustainability trade-off:} Minimizing IT energy need not minimize carbon, water or biodiversity impacts across deployment regions and times.
\item \key{Scope matters:} For a fixed workload and region/time, all four operational-only rankings follow configuration energy. Deployment or embodied impacts can change rankings.
\item \key{Outcome:} 50.2\% median reduction in worst-case regret vs.\ offline WaterWise in the 12-region evaluation (20 capacity seeds).
\end{summarysection}
\end{minipage}
\vspace{.18cm}
\begin{summarysection}{Method}
\item \key{Model:} Combine IT energy with regional environmental intensities to evaluate four operational impacts. Add embodied impacts for lifecycle comparisons.
\item \key{Compare:} Separate GPU/model configuration rankings from region/time deployment rankings. Identify when operational rankings agree and embodied impacts reverse them.
\item \key{Route:} Allocate workloads across regions to minimize the largest relative regret vs.\ each dimension's optimum, subject to regional capacity and serving constraints.
\end{summarysection}
\end{frame}

% Category A; sources are intentionally not shown on slides.
\categorycover{HBF System Characterization\\and KV-Centric Serving}{When does HBF improve LLM serving, and when does it become a bottleneck?}{Characterizing HBF; HBF Sucks?; Hot--Cold Tiering}
% Source: https://arxiv.org/abs/2609.39131v2
\begin{frame}[t]{Characterizing High Bandwidth Flash for LLM Serving}
\vspace{.02cm}
\begin{minipage}[t][3.25cm][t]{\linewidth}
\begin{summarysection}{Motivation}
\item \key{Capacity pressure:} Long contexts and agentic calls compete for HBM space for weights, active KV and reusable prefixes.
\item \key{Retention matters:} More HBF capacity helps only when placement and scheduling account for read costs, slow writes and finite endurance.
\item \key{Outcome:} Trace-driven simulation improves completion time and estimated write lifetime. Modeled energy savings exclude static power.
\end{summarysection}
\end{minipage}
\vspace{.18cm}
\begin{summarysection}{Method}
\item \key{Organize:} Model an HBM--HBF--host hierarchy and compare weight placement and KV residency under different memory organizations.
\item \key{Retain:} Track prefix reuse and cache movement. Buffer growing KV in HBM and preserve useful prefixes to avoid reloads or recomputation.
\item \key{Schedule:} Prioritize cache hits and limit admission with buffered cache-aware scheduling, balancing batch growth against future KV reuse.
\end{summarysection}
\end{frame}
% Source: https://arxiv.org/abs/2608.11668v4
\begin{frame}[t]{HBF Sucks? A Full-Stack Characterization of\\High-Bandwidth Flash for KV-Centric LLM Serving}
\vspace{.02cm}
\begin{minipage}[t][3.25cm][t]{\linewidth}
\begin{summarysection}{Motivation}
\item \key{Drop-in risk:} Faster flash does not ensure faster serving when HBF replaces SSD storage but also consumes near-tier package resources.
\item \key{Traffic mismatch:} Hot KV remains in the near tier, leaving HBF a cold, write-heavy stream with limited useful reads to repay writes.
\item \key{Outcome:} The modeled Mooncake-style substitution can worsen latency and goodput. This is a conditional deployment result, not a rejection of HBF.
\end{summarysection}
\end{minipage}
\vspace{.18cm}
\begin{summarysection}{Method}
\item \key{Replay:} Extend TokenSim with GPU/near-tier/HBF costs and replay production KV-offload traces using capacity-matched SSD baselines.
\item \key{Diagnose:} Test three conditions: reads on the critical path, enough useful reads per write, and bandwidth that can be sustained.
\item \key{Constrain:} Combine thermal throttling and write-endurance models to distinguish peak device bandwidth from practical serving benefits.
\end{summarysection}
\end{frame}
% Source: https://arxiv.org/abs/2609.25782v1
\begin{frame}[t]{Hot--Cold Tiering of HBM and High Bandwidth Flash\\for Agentic LLM Serving}
\vspace{.02cm}
\begin{minipage}[t][3.25cm][t]{\linewidth}
\begin{summarysection}{Motivation}
\item \key{Idle state:} Agentic sessions pause for tools or users while retaining large KV histories that exceed HBM capacity.
\item \key{Two access patterns:} Active KV is read every decode step; paused KV is read on resume. Serving all KV from HBF wastes read energy.
\item \key{Outcome:} Trace-based evaluation expands retained-session capacity while keeping active decoding in HBM and reducing flash read power.
\end{summarysection}
\end{minipage}
\vspace{.18cm}
\begin{summarysection}{Method}
\item \key{Classify:} Use session activity to separate a small hot KV working set from a large cold pool of paused-session state.
\item \key{Place:} Keep active KV in HBM. Evict cold KV to HBF under capacity pressure; immutable blocks with valid HBF copies need no rewrite.
\item \key{Resume:} Prefetch paused KV back into HBM through buffered inter-tier transfers before execution resumes, overlapping movement with decode.
\end{summarysection}
\end{frame}

% Category B; sources are intentionally not shown on slides.
\categorycover{HBM/HBF Architecture\\and Data-Path Co-Design}{How should GPU, HBM, and HBF be organized to improve memory efficiency?}{H3; FlashAccel; HBFlex; Beyond Capacity}
% H3 content pending access to the original paper.
\begin{frame}[c]{H3: Hybrid Architecture Using High Bandwidth Memory\\and High Bandwidth Flash for Cost-Efficient LLM Inference}
\centering
{\Large\color{Muted}Pending}
\end{frame}
% Source: https://arxiv.org/abs/2607.10186v2
\begin{frame}[t]{FlashAccel: Leveraging High-Bandwidth Flash (HBF)\\for High-Throughput LLM Inference}
\vspace{.02cm}
\begin{minipage}[t][3.25cm][t]{\linewidth}
\begin{summarysection}{Motivation}
\item \key{Capacity wall:} Large model weights and KV histories restrict batch size and prefix retention in HBM-based GPUs.
\item \key{Utilization gap:} HBF capacity alone cannot hide NAND latency or keep its many planes busy under changing inference accesses.
\item \key{Outcome:} The simulated hardware/software design improves throughput and energy efficiency over HBM-only serving under a latency constraint.
\end{summarysection}
\end{minipage}
\vspace{.18cm}
\begin{summarysection}{Method}
\item \key{Hide latency:} Integrate GPU, HBM and HBF with distributed SRAM buffers that prefetch data ahead of GPU consumption.
\item \key{Arrange data:} Lay out weights in access order and balance KV data across HBF planes to improve parallel reads and manage flash erasure.
\item \key{Manage:} Expose heterogeneous resources through an HBF-aware storage layer and programming model, coordinating persistent data and execution.
\end{summarysection}
\end{frame}
% Source: https://arxiv.org/abs/2609.18675v1
\begin{frame}[t]{HBFlex: A Flexible Memory System for Bridging Fine-Grained\\LLM States and Coarse-Grained HBF Parallel Execution}
\vspace{.02cm}
\begin{minipage}[t][3.25cm][t]{\linewidth}
\begin{summarysection}{Motivation}
\item \key{Package trade-off:} Full-HBF designs devote more package resources to capacity and parallel bandwidth than hybrid designs retaining HBM.
\item \key{Dynamic-state cost:} Fine-grained KV reads imbalance planes; incremental writes block reads; mixed KV lifetimes amplify garbage collection.
\item \key{Outcome:} Trace-driven simulations improve throughput over the evaluated H3 and FlashAccel baselines through coordinated KV management.
\end{summarysection}
\end{minipage}
\vspace{.18cm}
\begin{summarysection}{Method}
\item \key{Balance reads:} Distribute KV placement and attention accesses across planes, exposing parallelism and reducing overloaded-plane delays.
\item \key{Schedule writes:} Aggregate incremental KV updates and write them back during sufficiently long compute windows to reduce read interference.
\item \key{Reclaim carefully:} Pack KV with similar lifetimes into blocks and defer reclamation, reducing valid-page migration and extra flash writes.
\end{summarysection}
\end{frame}
% Source: https://drive.google.com/file/d/1IHfgK1Tj3ytcPLpDlG87RLXM27l9DEtv/view
\begin{frame}[t]{Beyond Capacity: Scalable MoE LLM Inference via\\High-Bandwidth Flash with Direct GPU and HBM Paths}
\vspace{.02cm}
\begin{minipage}[t][3.25cm][t]{\linewidth}
\begin{summarysection}{Motivation}
\item \key{Expert capacity:} MoE expert weights outgrow HBM even though each token uses only a small subset of experts.
\item \key{Delivery bottleneck:} A single HBM relay path can leave direct GPU--HBF bandwidth unused, while late expert selection exposes NAND latency.
\item \key{Outcome:} Event-driven serving simulation shows higher expert-delivery efficiency and throughput than either delivery route alone.
\end{summarysection}
\end{minipage}
\vspace{.18cm}
\begin{summarysection}{Method}
\item \key{Use two paths:} DASH assigns whole experts to direct GPU--HBF or HBF--HBM-base-die--GPU routes and transfers them concurrently.
\item \key{Look ahead:} Determine upcoming experts early enough to overlap HBF reads with preceding computation, reducing exposed read latency.
\item \key{Separate traffic:} Isolate immutable expert weights from mutable KV; buffer small decode writes and coordinate larger prefill transfers.
\end{summarysection}
\end{frame}

% Category C; sources are intentionally not shown on slides.
\categorycover{HBF Simulation\\and Power/Thermal Modeling}{How can HBF architectures be evaluated before physical hardware becomes available?}{HBFSim (real GPU); HBFSim (event-driven); HBM Power}
% Source: https://arxiv.org/abs/2609.09800v3
\begin{frame}[t]{HBFSim: Fast and Faithful Simulation of\\High-Bandwidth Flash Under Real GPU Execution}
\vspace{.02cm}
\begin{minipage}[t][3.25cm][t]{\linewidth}
\begin{summarysection}{Motivation}
\item \key{Evaluation gap:} Physical HBF is unavailable for validation; cycle-level GPU simulation is slow and fixed traces miss runtime changes.
\item \key{Core idea:} Run the real workload on a GPU and model only the program-visible HBF memory behavior.
\item \key{Outcome:} The paper demonstrates timing injection and logical-capacity experiments. Injection accuracy is not future-HBF prediction accuracy.
\end{summarysection}
\end{minipage}
\vspace{.18cm}
\begin{summarysection}{Method}
\item \key{Intercept:} Register modeled address ranges and instrument supported PTX memory accesses while real GPU kernels execute.
\item \key{Model service:} Timing-only mode retains GPU data; Capacity mode uses backing storage and GPU frames. MQSim supplies media timing.
\item \key{Check scope:} Public defaults use issue-point waits. Dependency waits are experimental; thermal/reliability modeling is optional and uncalibrated.
\end{summarysection}
\end{frame}
% Source: https://arxiv.org/abs/2610.23503v1
\begin{frame}[t]{HBFSim: A Simulator for Exploring HBF\\in Memory Hierarchies for LLMs}
\vspace{.02cm}
\begin{minipage}[t][3.25cm][t]{\linewidth}
\begin{summarysection}{Motivation}
\item \key{Design uncertainty:} HBM/HBF capacity, NAND latency and finite endurance make hierarchy choices difficult before hardware is available.
\item \key{Flash semantics:} Placement must account for page programming, out-of-place updates, garbage collection and data movement.
\item \key{Outcome:} The event-driven study finds active KV retention in HBM preferable to using HBM merely as an HBF cache in its evaluated workloads.
\end{summarysection}
\end{minipage}
\vspace{.18cm}
\begin{summarysection}{Method}
\item \key{Drive workloads:} Use a calibrated LLM serving frontend to generate weight, KV and scratch accesses with application dependencies.
\item \key{Simulate hierarchy:} Model HBM, HBF and host tiers with event-driven resource calendars, links and host-managed flash translation.
\item \key{Explore policies:} Vary placement and transfers; track NAND reads, programs, erases and write amplification to estimate serving time and endurance.
\end{summarysection}
\end{frame}
% Source: https://arxiv.org/abs/2610.14244v1
\begingroup
\setbeamerfont{frametitle}{size=\normalsize,series=\bfseries}
\begin{frame}[t]{Understanding the Power Consumption of Modern High Bandwidth Memory:\\Experimental Analysis and Modeling Using Real HBM2 DRAM Chips}
\vspace{.02cm}
\begin{minipage}[t][3.25cm][t]{\linewidth}
\begin{summarysection}{Motivation}
\item \key{Power uncertainty:} A fixed HBM power assumption misses differences across commands, banks and transferred data patterns.
\item \key{Measured variation:} Real HBM2 measurements show that internal structure and bus toggling materially change memory power.
\item \key{Outcome:} An empirical HBM model predicts measured power more accurately than evaluated prior models. This work characterizes HBM, not HBF.
\end{summarysection}
\end{minipage}
\vspace{.18cm}
\begin{summarysection}{Method}
\item \key{Measure:} Run controlled command and data-pattern experiments on 36 real HBM2 chips, separating structural and data-dependent effects.
\item \key{Model:} Extend DRAMPower using measured currents and switching behavior to account for bank organization and data transitions.
\item \key{Validate:} Compare predicted power with measurements and workload traces, including LLM decoding; separately check HBM3E GPU power.
\end{summarysection}
\end{frame}
\endgroup

% Category D; sources are intentionally not shown on slides.
\categorycover{HBF-Accelerated Retrieval\\and RAG Systems}{Can HBF accelerate retrieval workloads beyond LLM decoding?}{HAVEN: vector search and reranking}
% Source: https://drive.google.com/file/d/1B2Sy8GHSAgfoiHFMcmndiJHMQ1G9KODf/view
\begin{frame}[t]{HAVEN: High-Bandwidth Flash Augmented Vector Engine for\\Large-Scale Approximate Nearest-Neighbor Search Acceleration}
\vspace{.02cm}
\begin{minipage}[t][3.25cm][t]{\linewidth}
\begin{summarysection}{Motivation}
\item \key{Retrieval bottleneck:} Billion-scale RAG vector databases exceed GPU HBM. High-recall reranking must fetch full-precision vectors from DRAM or SSD.
\item \key{Locality opportunity:} On-package HBF can retain large read-mostly vector data near the GPU and reduce off-GPU transfers.
\item \key{Outcome:} Modeled IVF-PQ pipelines improve reranking throughput and latency over GPU--DRAM/SSD systems. The target is retrieval, not LLM decode.
\end{summarysection}
\end{minipage}
\vspace{.18cm}
\begin{summarysection}{Method}
\item \key{Store nearby:} Augment a GPU/HBM package with HBF holding full-precision vectors, avoiding host-memory round trips during reranking.
\item \key{Search in stages:} Use IVF-PQ to identify candidates, then a near-storage search unit to support exact-distance reranking close to vector data.
\item \key{Evaluate limits:} Model NAND subarrays, power-constrained HBF bandwidth and the retrieval pipeline to quantify data-movement and search benefits.
\end{summarysection}
\end{frame}

% Candidate research questions synthesized from the reviewed work; not proven novelty.
\begin{frame}[c]{What Remains Open in HBM/HBF Co-Design?}
\reviewpoint{Architecture trade-offs}{Under a common power and package budget, when do HBM-only, hybrid HBM/HBF and full-HBF designs offer the best capacity, bandwidth and serving efficiency?}
\reviewpoint{Power and thermal interaction}{How do static/dynamic power, throttling, retention maintenance and extra flash traffic change sustainable bandwidth and energy per request?}
\reviewpoint{Runtime and data placement}{Which tiers suit weights, active KV, paused KV, reusable prefixes and MoE experts? When do migration costs outweigh reuse benefits?}
\reviewpoint{Simulation fidelity}{Which conclusions survive changes in hardware assumptions, latency-injection semantics and thermal models without physical HBF validation?}
{\color{Muted}Candidate open questions from this review; no claim of first discovery.}
\end{frame}

\begingroup
\renewcommand{\reviewpoint}[2]{\par{\bfseries\color{Accent}#1}\par #2\par\smallskip}
\begin{frame}[t]{Proposed Research Directions and Next Steps}
\vspace{0cm}
{\bfseries\color{Accent}Power- and Thermal-Aware HBM/HBF Co-Design\\for Energy-Efficient LLM Serving}\par\medskip
\reviewpoint{1. Power-constrained architecture co-design}{Compare stack/capacity allocations in HBM-only, hybrid and full-HBF designs under common package limits. Minimize energy subject to serving SLOs.}
\reviewpoint{2. Thermal- and endurance-aware KV placement}{Separate active, paused and reusable KV. Compare adaptive and static placement for energy, temperature and flash writes.}
\reviewpoint{3. Simulator-based evaluation}{Reproduce both HBFSim baselines. Calibrate power/thermal models where justified and test sensitivity to uncertain hardware assumptions.}
{\bfseries\color{Accent}Next steps}\par\smallskip
Read 1--2 papers/week; reproduce baselines; test one hypothesis/iteration; refine scope with the advisor.
\par{\small\color{Muted}Candidate directions; novelty and thesis scope remain to be established.}
\end{frame}

\endgroup

\end{document}
